\section{Classical single-family results}\label{sup:classical-comparison}
For one scale, the resonant Bessel triple is classical. Mahler proved
independence of values of $J_0,J'_0$ and order derivatives of normalized
$J_\lambda$ at $\lambda=0$ \cite[Theorem~2, p.~167]{Mahler1968}. His
Theorem~3, pp.~169--170, rewrites these values with connection constants.
In particular, he proves that $\pi Y_0(2)/(2J_0(2))-\gamma$ is
transcendental \cite[Theorem~3 and its consequence, pp.~169--170]{Mahler1968}.
His $K_\lambda$ denotes normalized ordinary Bessel $J$, not modified $K$.
Belogrivov treated each
nonnegative integer order at nonzero algebraic points and the determinant
relation $z(u'_1u_0-u_1u'_0)+u_0^2=1$
\cite[Theorem~4, Lemmas~4--5, (36)]{Belogrivov1973}.
His coordinates are $u_{0,\lambda}(\zeta)=K_\lambda(\zeta)$ and
$u_{1,\lambda}(\zeta)=\partial_\lambda K_\lambda(\zeta)$, with the
same normalized ordinary Bessel convention. Thus
$u_{0,0}(icz)=I_c(z)$ and $u_{1,0}(icz)=-S_c(z)$.
His primes are derivatives in $\zeta$; substitution gives
$z(I_cW_c-V_cS_c)-I_c^2+1=0$.
His Kummer-pair theorems treat values of ${}_1F_1(\mu;\nu;z)$ and its
derivative at several points, without the logarithmic companion
\cite[Theorems~1--2]{Belogrivov1971}.
The one-point corollary over the full $\mathcal E$ is known from
Shidlovskii's simultaneous exponential-integral theorem
\cite[Chapter~7, \S1, Theorem~1, pp.~232--233]{Shidlovskii1989}
by the finite-rate reduction below.
Mahler's final integral family was announced without proof
\cite[\S23, pp.~172--173]{Mahler1968}, as Rivoal notes
\cite[\S7]{Rivoal2012}.
Put $G_0(z)=\int_0^\infty e^{-t}(t+z)^{-1}\,dt$ for
$z\in\C\setminus(-\infty,0]$, and use the principal branch of $\log z$
\cite[p.~1 of the author manuscript]{Rivoal2012}.
Rivoal proves his Theorem~2(ii) for every algebraic
$z\notin(-\infty,0]$ \cite[Theorem~2(ii), \S\S5 and~7]{Rivoal2012}:
$\trdeg_{\Q}\Q(\gamma+\log z,e^z,G_0(z))\ge2$.
At $z=1$, $G_0(1)=\delta=eE_1(1)$, so this gives
$\trdeg_{\Q}\Q(e,\gamma,\delta)\ge2$.
The full $\mathcal E$-base conclusion is attributed to Shidlovskii's
simultaneous theorem, without adding an all-rate quantifier to Rivoal's statement.
Main-text Theorem~\ref{thm:connection} systematically assembles
multi-node mixed families over the selected regular-coordinate field $B$.
The source comparison for that scope remains partial.

\subsection*{Finite-rate reduction for the known one-point consequence}
Shidlovskii defines $\omega_0(z)=e^z$ and
$\omega_k(z)=1+\sum_{n\ge1}z^n/(n!n^k)$ for $k\ge1$.
Thus $\omega_1(-x)=1-\Ein(x)$.
For each integer $m\ge1$ and each finite set of algebraic nodes
$\xi_1,\ldots,\xi_r$ linearly independent over $\Q$, his theorem gives
simultaneous algebraic independence of $\omega_k(\xi_i)$,
$0\le k\le m$, $1\le i\le r$. The reduction below uses $m=1$.
Fix algebraic $x\ne0$. A polynomial relation for $\Ein(x)$ over
$\mathcal E$ has finitely many coefficients, each algebraic over a
field generated by finitely many $e^{\beta_j}$ with algebraic $\beta_j$.
Extend $\eta_1=-x$ to a $\Q$-basis $\eta_1,\ldots,\eta_r$ of the
span of $-x$ and these rates.
Clearing rational denominators makes every $e^{\beta_j}$ algebraic
over $K=\Qbar(e^{\eta_1},\ldots,e^{\eta_r})$.
The theorem with $m=1$ makes $\omega_1(-x)$ transcendental over $K$.
It remains transcendental over any algebraic extension containing the
finitely many coefficients, excluding the proposed relation.
At $x=1$, $\Ein(1)=\gamma+E_1(1)$ and $e\in\mathcal E^\times$,
giving the equivalent alternatives with $E_1(1)$ and $\delta=eE_1(1)$.
This reduction retains $-x$ as a basis member. It does not retain
arbitrary merely distinct nodes in a $\Q$-basis or adjoin regular
Bessel/Kummer pairs; the mixed conclusion over $B$ uses
main-text Theorem~\ref{thm:values}.

Powers treats odd and even multisections of iterated exponential integrals.
He proves independence over the field generated by their level-zero functions.
At $1$, he proves independence over the corresponding value field \cite[p.~13]{Powers2026}.
Belogrivov also treated one-node shifted-denominator hypergeometric
families and argument derivatives \cite{Belogrivov1970}.

\subsection*{Homogeneous and integration antecedents}
Ostrowski makes primitive dependence explicit for rational-exponential
derivatives \cite[pp.~315--316]{Ostrowski1946}; Kolchin gives the general
differential-algebraic setting \cite{Kolchin1973}.
Beukers, Brownawell and Heckman's 1988 paper
\cite{BBH1988} is an antecedent for the differential-Galois method.
BBH's homogeneous framework supplies its normality characterization
\cite[Theorem~2.2, p.~288]{BBH1988}.
Its common-power descent assumes irreducible Siegel-normal systems
\cite[Proposition~2.4, pp.~291--292]{BBH1988}.
Gorelov invokes these results in his homogeneous reduction
\cite[Lemma~13, p.~7]{Gorelov2021}.
He identifies BBH's omitted Bessel--Kummer linkage
\cite[Example~1, p.~3; Lemma~14, p.~7]{Gorelov2021}.
The second BBH result cited there is Proposition~2.4, not Theorem~2.4.
His criteria characterize independence of fundamental-matrix entries with exponentials
and of the values $F^{(i)}(\alpha)$ \cite[Theorems~1--2, pp.~4--5]{Gorelov2021}.
Appendix~\ref{app:gorelov} checks the hypotheses for the homogeneous rank-two inputs used here.
BBH Theorem~2.2(iv) requires an extension by scalars of $\SL_2$.
This is the rank-two group condition supplied by main-text Lemma~\ref{lem:groups}.
It includes the homogeneous logarithmic solution $J=S-I L$ or $J=\Sigma-FL$.
No normality assertion applies to the enlarged extension with $L$ and the E-parts separately.

Gorelov's inhomogeneous generalized hypergeometric equation has scalar forcing
$\prod_{j=1}^q(\lambda_j-1)$ \cite[p.~663]{Gorelov2016}.
His inhomogeneous hypergeometric criteria require homogeneous-group and parameter hypotheses
\cite[Theorems~1--3, pp.~664--665]{Gorelov2016}.
For $q>l$, his function criteria work over homogeneous fundamental-matrix fields.
They require a homogeneous group containing $\SL_q$
\cite[Theorems~1--2, p.~664]{Gorelov2016}.
Dependence occurs when a denominator parameter belongs to $\mathbb N$;
for a collection, equivalent parameter vectors at equal scales also cause dependence.
His value theorem requires rational parameters, odd $q\ge3$, irreducibility exclusions and nonzero algebraic arguments.
It also imposes four explicit collision conditions
\cite[Theorem~3, pp.~664--665]{Gorelov2016}.
These hypotheses are not conditions for our normalized logarithmic E-parts.
In our resonant encoding, $\lambda=(1,1)$ makes that scalar forcing vanish.
The E-parts instead have forcing $2I'$ or $2F'-sF$.
Their name alone does not identify the two inhomogeneous problems.

Gorelov's primitive dependence criterion is a constant-linear-combination criterion
\cite[Lemma~6, p.~670]{Gorelov2016}.
His integration result adjoins primitives of $z^\delta e^{\omega z}$
over the specified homogeneous fields
\cite[Lemma~7, pp.~670--672]{Gorelov2016}.
Here $\delta$ is his integration exponent, distinct from the Gompertz constant above.
The nonzero rates are distinct, $\delta\notin\Z_{\ge0}$,
and same-rate exponents do not differ by integers.
Independence of his fundamental-matrix, exponential and power functions is assumed.
The case $\delta=-1$ includes first-level exponential integrals.
Lemma~8 concerns integration involving homogeneous solution cofactors.
These are integration antecedents.
They do not state the full finite-pole residue algorithm or its affine normalization constants.

BBH already discusses successive exponential integrals as a Liouville extension
\cite[p.~281]{BBH1988}.
Its successive series are $\sum_{n\ge1}z^n/(n^j n!)$ through finite depth.
They arise from successive inhomogeneous first-order equations, outside Siegel normality.
This anticipates tower integration, without giving the present multi-node tower ideal or its degree.

\subsection*{Gorelov's normalized integral family}
For $\lambda\in\C\setminus\Z_{<0}$, write
$\phi_\lambda(z)={}_1F_1(1;\lambda+1;z)=\sum_{n\ge0}z^n/(\lambda+1)_n$,
so $\phi_\lambda(0)=1$ \cite[p.~384]{Gorelov2025}.
Gorelov's later paper treats the specific family
$V_{\lambda,\nu,\alpha}(z)=\allowbreak\nu z^{-\nu}e^{\alpha z}\allowbreak\int_0^z t^{\nu-1}e^{-\alpha t}\phi_\lambda(t)\,dt$
\cite[pp.~383--384]{Gorelov2025}.
Here $\alpha\in\C$ and $\lambda,\nu\in\C\setminus\Z_{<0}$.
Explicitly, if $e^{-\alpha z}\phi_\lambda(z)=\sum_{n\ge0}a_nz^n$, then
$V_{\lambda,\nu,\alpha}(z)=e^{\alpha z}
\bigl(a_0+\sum_{n\ge1}\nu a_nz^n/(\nu+n)\bigr)$.
This is Shidlovskii's normalized primitive for $\nu\notin\Z_{<0}$
\cite[p.~384]{Gorelov2025}; in particular, $V_{\lambda,0,\alpha}=e^{\alpha z}$.
The integral need not converge at zero for every allowed complex parameter.
This family has $V_{\lambda,\nu,\alpha}(0)=1$.
It satisfies a second-order inhomogeneous equation and
$V'=(\alpha-\nu/z)V+(\nu/z)\phi_\lambda$
\cite[p.~384; (17), p.~385]{Gorelov2025}.
Theorem~1 classifies its Laurent-polynomial representations in exponentials and $\phi$ functions
for the stated complex parameters. For members of this family that are E-functions,
Theorem~2 classifies algebraic dependence of $V$ and $V'$ over $\C(z)$
\cite[pp.~384--385]{Gorelov2025}.
Lemma~3 gives a second-order inhomogeneous characterization of E-functions
\cite[p.~386]{Gorelov2025}.
Its value discussion concerns that family
\cite[Conclusion~1, p.~390]{Gorelov2025}.
Collections $V_{\lambda,\nu,\alpha}(\xi z)$ with varying parameters and scales are explicitly deferred
\cite[\S4, item~2 (Conclusion~2), p.~390]{Gorelov2025}.

These works supply homogeneous, integration and value antecedents.
They do not state the joint multi-scale normalized ideal with shared logarithmic E-parts,
finite-pole residue relations, higher Ein towers and fixed additive constants.
They also do not give the present actual connection-constant cancellation algebra.
No global novelty conclusion follows from this comparison.

\subsection*{Earlier value inputs}
Earlier value theorems have narrower inputs.
BBH treats rational admissible homogeneous hypergeometric families
\cite[Corollary~4.6, pp.~300--301]{BBH1988}.
Gorelov imposes rational-parameter and collision conditions
\cite[Theorem~3, pp.~664--665]{Gorelov2016}.
His later result concerns one integral family, with Theorem~2 restricted to E-functions
\cite[Theorem~2, p.~385; Conclusion~1, p.~390]{Gorelov2025}.
These statements do not supply the present specialized normalized ideal.
